\section{偏序与可比较性}
\label{sec:01-Mathematical-Language/02-Sets-Relations-Functions/03}

选型会上摆着三个候选模型。\(A\) 的准确率 \(0.91\)，推理延迟 120 毫秒；\(B\) 的准确率 \(0.88\)，延迟 40 毫秒；\(C\) 的准确率 \(0.87\)，延迟 150 毫秒。有人问：哪个最好？

\(C\) 好办——它在两项指标上都不如 \(A\)，可以当场淘汰，谁也不会为它辩护。\(A\) 与 \(B\) 之间却给不出答案，除非有人先说清楚一毫秒延迟值多少个百分点的准确率。这不是数据不够，也不是没算完：把两项指标当作同等重要，\(A\) 与 \(B\) 之间就是没有大小可言。

\begin{center}
\begin{tikzpicture}[font=\small,>=stealth]
  \fill[black!8] (3.6,0) rectangle (6.3,2.75);
  \draw[->,black!55] (0,0) -- (6.7,0) node[right,black] {延迟};
  \draw[->,black!55] (0,0) -- (0,3.4) node[above,black] {准确率};
  \fill (3.6,2.75) circle (2pt);
  \fill (1.2,1.25) circle (2pt);
  \fill (4.5,0.75) circle (2pt);
  \node[anchor=south west] at (3.68,2.8) {\(A\)};
  \node[anchor=south west] at (1.28,1.3) {\(B\)};
  \node[anchor=south west] at (4.58,0.8) {\(C\)};
  \node[black!60] at (5.35,2.15) {被 \(A\) 支配};
\end{tikzpicture}

\emph{延迟越小越好、准确率越大越好，于是阴影区里的点在两项上都不优于 \(A\)。\(C\) 落在里面，可以淘汰；\(B\) 既不在 \(A\) 的阴影里，也没把 \(A\) 装进自己的阴影——两者不可比。}
\end{center}

比较关系只覆盖一部分对象对，这不是个例。集合 \(\{1\}\) 与 \(\{2\}\) 互不包含，正整数 2 与 3 互不整除，两个训练任务之间可能根本没有先后要求。实数上那种``任意两个都能比''的顺序，反倒是特例。

\subsection{偏序：三条要求，一条比一条挑剔}

\begin{definition}
集合 \(P\) 上的关系 \(\preceq\) 若满足自反性（对每个 \(x\) 有 \(x\preceq x\)）、反对称性（\(x\preceq y\) 且 \(y\preceq x\) 蕴含 \(x=y\)）和传递性（\(x\preceq y\) 且 \(y\preceq z\) 蕴含 \(x\preceq z\)），就称为偏序，\((P,\preceq)\) 称为偏序集。
\end{definition}

自反和传递与上一节的等价关系共用，分岔发生在第二条：等价关系要对称，偏序要反对称。对称让关系横着摊开、把对象聚成类，反对称禁止不同对象互相排在对方前面，从而把对象竖着码成层次。注意反对称并不是``不对称''——自反性明明要求 \(x\preceq x\)，所以自己指向自己是允许的，被禁止的只是两个不同对象之间的双向箭头。

三条里哪条最容易失守？在模型选型这个例子里恰好是反对称。若两个不同的模型在所有记录的指标上完全一致，它们互相支配却不是同一个模型，反对称性当场作废。这时候得到的只是预序（preorder）。修补方式有两条：要么承认它是预序、放弃唯一性的结论，要么先按``指标向量相同''这个等价关系取商，在商集上再谈偏序。上一节造有理数用的是同一招——先把该视为相同的东西合并掉，再在类上定义结构。

\subsection{两个立刻能上手的例子}

幂集上的包含关系是最标准的偏序。取 \(S=\{1,2,3\}\)，\(\mathcal P(S)\) 的八个元素里，\(\{1\}\subseteq\{1,2\}\) 成立，而 \(\{1\}\) 与 \(\{2\}\) 互不包含。反对称性直接由外延性给出：双向包含即相等。

正整数上的整除关系稍微费点事。自反取 \(k=1\)；传递是把两次倍数关系乘起来。反对称的骨架是：由 \(b=k_{1}a\) 与 \(a=k_{2}b\) 得 \(a=k_{1}k_{2}a\)，在正整数里 \(a>0\)，故 \(k_{1}k_{2}=1\)，而正整数相乘得 1 只能两个都是 1，于是 \(a=b\)。

请留意这一步用在哪里：整个论证的重量全压在``正整数''上。把论域放宽到全体整数，反例立刻出现——\(1\mid -1\) 且 \(-1\mid 1\)，但 \(1\neq -1\)。同一个符号 \(a\mid b\)，换个地盘就不再是偏序。判断一条关系的性质时，光看符号不够，得看它踩在哪个集合上。

如果偏序集里任意两个元素都可比，就称为全序。实数上的 \(\le\) 是全序，包含与整除都不是。日常直觉几乎全部来自全序，而实际建模中遇到的多半是偏序：任务依赖、集合族的包含、多目标的支配、函数的逐点比较，无一是全序。

\subsection{最大、极大、上界：一个例子说清三个词}

仍取正整数与整除关系，看子集 \(A=\{2,3,6\}\)。

6 是 \(A\) 的最大元：它自己在 \(A\) 里，且 \(A\) 的每个元素都整除它。最大元若存在必唯一——两个最大元互相排在对方前面，反对称性一压，它们是同一个。

换成 \(B=\{2,3\}\)，最大元就没有了：2 不被 3 整除，3 也不被 2 整除，\(B\) 里没有谁能压住所有人。但 2 是极大元，因为 \(B\) 里没有比它严格更大的；同理 3 也是极大元。极大元只要求``没人在它上面''，不要求``它在所有人上面''，因此可以有好几个互不可比的极大元并存。

上界则放弃了``必须在子集内''这一条。对 \(B=\{2,3\}\)，6 是上界，12 和 18 也是；所有上界中最小的那个是 6，称为上确界。6 并不属于 \(B\)，却是所有盖住 \(B\) 的元素里最紧的一个。上确界存不存在、是否唯一，是后面很多分析学结论的支点，\hyperref[ch:097]{``实数为什么需要完备''}讨论的正是实数系统在这一点上的特殊之处。

这里最容易出岔子的是把偏序的``大''读成数值的``大''。在整除偏序里，6 比 2 和 3 都大，因为大意味着被整除；\(\{1,2\}\) 在包含偏序里比 \(\{1\}\) 大，因为大意味着包含得更多。每换一种偏序，都要重新校准``大''指的是什么。

\subsection{Hasse 图：把不可比画出来}

有限偏序集可以画成 Hasse 图。先记 \(x\prec y\) 表示 \(x\preceq y\) 且 \(x\neq y\)。若 \(x\prec y\) 且中间不存在 \(z\) 满足 \(x\prec z\prec y\)，就把 \(y\) 画在 \(x\) 上方并连一条线。自反的自环省略，能由传递性推出的长边也省略，只留最近的上层邻居。图里不画箭头，上下位置已经交代了方向。

用一条数据流水线做例子：清洗原始数据之后，生成模型特征与生成质量报告各自独立开工，两者都完成才能最终发布。

\begin{center}
\begin{tikzpicture}[
  every node/.style={draw, rounded corners=2pt, fill=blue!6,
    minimum width=3.2cm, minimum height=0.8cm, align=center},
  every path/.style={draw, line width=0.7pt}
]
  \node (clean) at (0,0) {清洗原始数据};
  \node (feature) at (-2.4,1.7) {生成模型特征};
  \node (report) at (2.4,1.7) {生成质量报告};
  \node (publish) at (0,3.4) {最终发布};

  \draw (clean) -- (feature);
  \draw (clean) -- (report);
  \draw (feature) -- (publish);
  \draw (report) -- (publish);
\end{tikzpicture}

\emph{任务依赖的 Hasse 图：位置越高，任务越晚；左右两支之间没有任何向上的通路，它们不可比。}
\end{center}

从``清洗原始数据''沿任意向上的路都能走到``最终发布''，所以清洗必然早于发布，但图上这两点之间没有直接连线——这条关系已由中间节点和传递性给出，画出来是冗余。左右两支之间上下都走不通，它们不可比。

\subsection{不可比不是未决定}

不可比在调度里是有价值的信息：它意味着两个任务可以并行。把它们强行排成先后，等于凭空给系统加了一条不存在的约束，本可以并发的工作被串行化，而任何监控都不会报错。

反方向的错误更隐蔽。两个任务在业务上明明有先后——一个准备数据，另一个用这份数据训练——但你在依赖定义里没写这条边。Hasse 图会诚实地反映你写下的定义，不可比就画成不可比，它不会替你脑补遗漏。图上少一条线，运行时就少一条约束，而事故往往要等到某次调度恰好把顺序颠倒时才暴露。

这类结构在现实里比树常见得多。幂集里 \(\{1,2\}\) 与 \(\{2,3\}\) 不可比，整除关系里 4 与 6 不可比，多目标评估里 \(A\) 与 \(B\) 不可比。这些空缺不是建模的粗糙，而是结构本身的形状，偏序给了描述这种形状的语言。

至此，同一集合上的关系已经分出两条路：对称的那条把对象归类，反对称的那条把对象分层。下一节转向两个集合之间的关系，并且只保留其中最受约束的一类——每个输入恰好配一个输出。约束一紧，就可以追问一件前面无法追问的事：这次对应能不能反着走回去。
